\documentclass[12pt,a4paper]{article}

\usepackage[utf8]{inputenc}
\usepackage[english]{babel}
\usepackage{geometry}
\usepackage{booktabs}
\usepackage{array}
\usepackage{longtable}
\usepackage{multirow}
\usepackage{multicol}
\usepackage{xcolor}
\usepackage{graphicx}
\graphicspath{{/home/pablo/Desktop/master/tfm/figures/}}
\usepackage{amsmath}
\usepackage{natbib}


\newcommand{\mnras}{MNRAS}
\newcommand{\apj}{ApJ}
\newcommand{\apjl}{ApJL}
\newcommand{\apjs}{ApJS}
\newcommand{\aap}{A\&A}
\newcommand{\aj}{AJ}
\newcommand{\araa}{ARA\&A}
\newcommand{\pasp}{PASP}
\newcommand{\pasj}{PASJ}
\newcommand{\nat}{Nature}
\newcommand{\sci}{Science}

\usepackage{float}
\usepackage{caption}
\usepackage{cuted}

\geometry{
    top=1.5cm,
    bottom=1.5cm,
    left=1.5cm,
    right=1.5cm
}

\usepackage[colorlinks=true, linkcolor=blue, citecolor=blue, urlcolor=blue]{hyperref}

\begin{document}

\section{Introduction}


The Cosmic Microwave Background (CMB) provides a unique window to study the early stages of the Universe and its structure formation. In particular, measurements of CMB polarization provide a powerful probe of primordial gravitational waves generated during the inflationary epoch. These gravitational waves imprint a distinctive pattern in the CMB polarization, known as B-modes.

Detecting primordial the B-mode would constitute direct evidence for inflationary gravitational waves, enabling a measurement of the tensor-to-scalar ratio, and thereby constraining the energy scale of inflation. Such a discovery would open an unprecedented observational window onto high-energy physics, far beyond the reach of terrestrial experiments.

However, its detection does not depend solely on instrumental sensitivity. A successful discovery critically relies on accurate component separation, since the cosmological signal is contaminated by polarized Galactic emission with significantly higher amplitude. 


\subsection{Galactic foregrounds and data}

In this figure, you can see the frequency dependence of the CMB in brightness temperature. 


\section{Methodology}

\subsection{Angular power spectra (APS)}

\subsection{Noise estimation}

\subsection{Window functions and CMB subtraction}

\subsection{Error estimation}

\section*{Foreground model}

\subsection{Synchrotron emission}

\subsection{Thermal dust emission}

\subsection{Synchrotron-dust spatial correlation}

\subsection{Bayesian inferenced and MCMC}

\subsection{Masks and summary}


\section{Results}

\subsection{Synchrotron amplitude ratio}

\subsection{Dust amplitude ratio}

\subsection{Synchrotron spectral index}

\subsection{Dust spectral index}

\subsection{Synchrotron-dust correlation coefficient}

\subsection{Synchrotron slope}

\subsection{Dust slope}

\section{Conclusions}

\end{document}